\chapter{What We Have Learned}

\section*{학습 목표}

이 장은 책 전체의 기술 내용을 다시 압축한다. 목적은 감상적 결론을 쓰는 것이 아니라, Real VLA를 읽고 설계할 때 반복해서 써야 하는 system-level vocabulary를 정리하는 것이다. 앞 장들에서 배운 핵심은 단순하다. VLA는 더 큰 VLM이 아니라, language-conditioned perception과 learned action을 state estimation, control, planning, data, evaluation, safety assurance에 연결하는 책임 구조이다.

\begin{enumerate}
  \item Part I--V가 어떤 순서로 Real VLA system claim을 닫는지 설명한다.
  \item $\obs_t,\state_t,\bel_t,\act_t,\ctrlcmd_t,\conf_t$의 역할을 다시 정리한다.
  \item state에서 action, control, data, evaluation, safety로 이어지는 연결 구조를 정식화한다.
  \item mug/rack/glass running example을 책 전체의 stack 관점에서 다시 해석한다.
  \item VLA 논문과 시스템을 읽을 때 반드시 확인해야 할 질문을 checklist로 정리한다.
\end{enumerate}

\section{The thesis revisited}

이 책의 첫 명제는 다음과 같았다.

\begin{definitionbox}{Real VLA thesis}
Real VLA는 특정 모델 이름이 아니다. Real VLA는 language-conditioned perception과 learned action을 실제 로봇의 물리 실행, monitoring, evaluation, safety assurance로 연결하는 책임 구조이다.
\end{definitionbox}

책 전체를 지나온 뒤 이 명제는 다음 식으로 다시 쓸 수 있다.
\begin{equation}
  \mathcal{R}
  =
  (
  \mathcal{O},
  \mathcal{B},
  \mathcal{A},
  \mathcal{U},
  \mathcal{D},
  \mathcal{E},
  \mathcal{S}
  ),
  \label{eq:ch21-real-vla-contract}
\end{equation}
여기서 $\mathcal{O}$는 observation interface, $\mathcal{B}$는 state/belief representation, $\mathcal{A}$는 policy-level action space, $\mathcal{U}$는 controller command space, $\mathcal{D}$는 data engine, $\mathcal{E}$는 evaluation protocol, $\mathcal{S}$는 safety assurance layer이다.

따라서 VLA system claim은 다음 질문으로 환원된다.
\begin{equation}
  \mathrm{Claim}
  =
  f(
  \mathrm{model},
  \mathrm{interface},
  \mathrm{data},
  \mathrm{control},
  \mathrm{evidence},
  \mathrm{safety}
  ).
  \label{eq:ch21-system-claim}
\end{equation}
model 항이 강해도 나머지 항이 비어 있으면 claim은 닫히지 않는다. 반대로 backbone이 보통 수준이어도 action interface, data record, evaluation protocol, safety gate가 명확하면 system contribution은 강할 수 있다.

\section{Part I: robot substrate}

Part I은 Real VLA가 올라가는 물리적 기반을 고정했다. 이 부분의 결론은 VLA policy가 image와 language만으로 로봇을 움직이는 것이 아니라는 점이다. 실제 실행은 state, dynamics, controller, contact, planning 위에서 일어난다.

\begin{equation}
  \bel_t
  =
  P(\state_t \mid \obs_{\leq t},\ctrlcmd_{<t}),
  \qquad
  \act_t
  \sim
  \pi_{\theta}(\cdot \mid \obs_{\leq t},\bel_t,\task).
  \label{eq:ch21-belief-policy}
\end{equation}

여기서 $\obs_t$는 sensor input이고, $\state_t$는 물리 상태이며, $\bel_t$는 uncertainty를 포함한 internal estimate이다. VLA가 직접 보는 것은 대개 $\state_t$가 아니라 $\obs_t$와 그로부터 추정된 representation이다. 이 차이를 무시하면 perception success와 robot success를 혼동한다.

control interface는 다음 mapping이다.
\begin{equation}
  \ctrlcmd_t
  =
  \controller(\act_t,\hat{\state}_t,C_t),
  \label{eq:ch21-action-control}
\end{equation}
여기서 $C_t$는 collision, joint limit, force limit, workspace, human proximity 같은 constraint이다. 좋은 action이라도 controller가 추적할 수 없거나 constraint를 위반하면 robot action이 아니다.

\begin{table}[h]
  \centering
  \caption{Part I에서 고정한 핵심 구분}
  \begin{tabularx}{\linewidth}{p{0.22\linewidth}X X}
    \toprule
    구분 & 의미 & VLA에서의 위험 \\
    \midrule
    Observation vs state & sensor input과 실제 물리 상태의 차이 & 보이는 것과 가능한 것을 혼동 \\
    Policy action vs control command & learned policy output과 actuator/controller command의 차이 & $\act_t$를 곧바로 $\ctrlcmd_t$로 오해 \\
    Geometry vs contact & 위치 관계와 접촉/힘 제약의 차이 & grasp, insertion, pushing 실패를 semantic error로 오해 \\
    Plan vs executable trajectory & plausible subgoal과 feasible motion의 차이 & LLM plan을 robot plan으로 과대해석 \\
    \bottomrule
  \end{tabularx}
  \label{tab:ch21-part1-summary}
\end{table}

\section{Part II: learning foundations}

Part II는 imitation learning, reinforcement learning, visuomotor policy를 통해 VLA policy의 직접 조상을 정리했다. 핵심은 VLA가 갑자기 새로 생긴 문제가 아니라, robot learning의 오래된 실패 모드가 language와 foundation model 위에서 다시 나타난다는 것이다.

behavior cloning은 다음 손실을 줄인다.
\begin{equation}
  \mathcal{L}_{BC}
  =
  \sum_{t}
  \ell(
  \pi_{\theta}(\obs_{\leq t},\task),
  \act_t^{E}
  ).
  \label{eq:ch21-bc-review}
\end{equation}
그러나 rollout에서는 policy가 만든 action이 다음 observation distribution을 바꾼다.
\begin{equation}
  \obs_{t+1}
  \sim
  P(\cdot \mid \state_t,\act_t),
  \qquad
  \act_t \sim \pi_{\theta}.
  \label{eq:ch21-rollout-shift}
\end{equation}
이것이 covariate shift이다. VLA가 language instruction을 이해하더라도, 실패 직전 state와 recovery action을 보지 못했다면 실제 rollout은 무너질 수 있다.

reinforcement learning과 optimal control은 reward, cost, constraint를 명시한다.
\begin{equation}
  J(\pi)
  =
  \mathbb{E}_{\pi}
  \left[
  \sum_{t=0}^{T}
  r(\state_t,\act_t,\task)
  -
  \lambda c(\state_t,\act_t)
  \right].
  \label{eq:ch21-rl-review}
\end{equation}
VLA 시대에도 이 식은 사라지지 않는다. 단지 reward가 language-conditioned task progress, physical safety, intervention, recovery, latency를 함께 포함해야 한다.

\section{Part III: from planner to VLA}

Part III는 VLM/LLM, planner+skills, robotics transformer, action representation, data engine, adaptation을 연결했다. 여기서 책의 중심 장은 action representation이다. VLA의 진짜 interface는 prompt가 아니라 $\act_t$이다.

\begin{equation}
  \act_t
  \in
  \{
  \mathrm{token},
  \mathrm{vector},
  \mathrm{waypoint},
  \mathrm{delta\ pose},
  \mathrm{chunk},
  \mathrm{distribution},
  \mathrm{whole\ body}
  \}.
  \label{eq:ch21-action-spectrum}
\end{equation}

각 action type은 downstream contract를 바꾼다.
\begin{equation}
  \mathrm{ActionContract}
  =
  (
  \mathcal{A},
  f_{VLA},
  f_{ctrl},
  \controller,
  \mathcal{C}_{safe},
  \mathrm{interrupt}
  ).
  \label{eq:ch21-action-contract-review}
\end{equation}
token action은 language-model training pipeline과 잘 맞지만 quantization과 decoding latency 문제가 있다. continuous action은 controller와 연결하기 쉽지만 normalization과 frame convention이 중요하다. chunk와 diffusion/flow action은 temporal consistency와 multimodality를 주지만 interruption과 safety projection이 어려워진다.

\begin{table}[h]
  \centering
  \caption{Action representation이 바꾸는 system contract}
  \begin{tabularx}{\linewidth}{p{0.24\linewidth}X X}
    \toprule
    Action type & 얻는 것 & 반드시 확인할 것 \\
    \midrule
    Token & sequence modeling과 semantic transfer & quantization, invalid token, decoder latency \\
    Delta pose & controller 연결과 빠른 servo & frame, scale, saturation, collision projection \\
    Waypoint & planning/control integration & IK, feasibility, local planner assumption \\
    Chunk & temporal consistency와 bimanual coordination & horizon, stale action, interruption rule \\
    Diffusion/flow & multimodal action distribution & sampling budget, safety validation, smoothness \\
    Whole-body & humanoid task coupling 표현 & balance, contact, self-collision, fall recovery \\
    \bottomrule
  \end{tabularx}
  \label{tab:ch21-action-contract-summary}
\end{table}

\section{Part IV: deployment evidence}

Part IV는 evaluation, hierarchy, real-time runtime, safety assurance를 다루었다. 이 부분의 핵심은 capability claim과 deployment evidence를 분리하는 것이다. benchmark success는 필요하지만 충분하지 않다.

\begin{equation}
  \mathcal{M}_{deploy}
  =
  (
  R_{success},
  R_{safe},
  T_{p95},
  E_{track},
  N_{intervention},
  R_{recovery}
  ).
  \label{eq:ch21-deploy-metrics}
\end{equation}
여기서 $R_{success}$만 있으면 성공률이다. $R_{safe}$, latency, tracking error, intervention, recovery가 함께 있어야 robot capability claim에 가까워진다.

hierarchical VLA는 이 claim을 layer별로 분해한다.
\begin{equation}
  \mathrm{Decision}
  =
  (
  d^{reason},
  d^{skill},
  \act_t,
  \ctrlcmd_t,
  d^{monitor}
  ).
  \label{eq:ch21-hierarchy-decision}
\end{equation}
어느 layer가 무엇을 결정했는지 기록하지 않으면 실패 원인을 알 수 없다. 상위 reasoning이 틀렸는지, skill selection이 틀렸는지, action representation이 틀렸는지, controller tracking이 틀렸는지, monitor가 늦었는지 구분해야 한다.

\section{Part V: beyond manipulation}

Part V는 humanoid와 future direction으로 범위를 넓힌다. humanoid VLA는 arm-only VLA의 단순 고차원 버전이 아니다. base, torso, head, arms, hands, legs, gaze, balance, locomotion, contact가 하나의 closed-loop system 안에서 결합된다.

\begin{equation}
  \act_t^{hum}
  =
  (
  \act_t^{base},
  \act_t^{torso},
  \act_t^{head},
  \act_t^{arms},
  \act_t^{hands},
  \act_t^{legs}
  ).
  \label{eq:ch21-humanoid-review}
\end{equation}
중요한 것은 dimension이 아니라 coupling이다. arm target은 balance를 바꾸고, foot placement는 reachability를 바꾸며, gaze는 perception quality를 바꾼다. 따라서 humanoid VLA claim에는 whole-body controller와 safety envelope이 반드시 포함되어야 한다.

Future Directions는 이 복습 위에서 읽어야 한다. world model, self-improving data loop, online adaptation, multi-embodiment transfer, safety assurance는 새로운 유행어가 아니다. 지금까지 정리한 interface와 evidence requirement가 커진 결과이다.

\section{Notation recap}

\begin{table}[h]
  \centering
  \caption{이 책에서 반복해서 쓰는 핵심 notation}
  \begin{tabularx}{\linewidth}{p{0.16\linewidth}X X}
    \toprule
    기호 & 의미 & 확인해야 할 질문 \\
    \midrule
    $\obs_t$ & observation, sensor input, language context & 무엇을 실제로 보았는가? calibration과 delay는? \\
    $\state_t$ & physical state & 무엇이 관측되지 않았고 추정되어야 하는가? \\
    $\bel_t$ & uncertainty를 포함한 belief & uncertainty가 action과 safety에 반영되는가? \\
    $\conf_t$ & robot configuration 또는 generalized coordinate & reachability, joint limit, kinematic feasibility는? \\
    $\act_t$ & policy-level action & token, pose, delta, chunk, distribution 중 무엇인가? \\
    $\ctrlcmd_t$ & controller/driver command & 어떤 controller가 어떤 rate로 실행하는가? \\
    $C_t$ & constraint set & collision, force, human proximity, workspace가 hard constraint인가? \\
    $m_i$ & data/embodiment metadata & robot, camera, controller, normalization이 기록되는가? \\
    \bottomrule
  \end{tabularx}
  \label{tab:ch21-notation-recap}
\end{table}

\section{The system loop}

책 전체를 하나의 loop로 쓰면 다음과 같다.
\begin{equation}
  \begin{aligned}
  \bel_t &= \mathrm{Estimate}(\obs_{\leq t},\ctrlcmd_{<t}),\\
  \act_t &\sim \pi_{\theta}(\cdot \mid \obs_{\leq t},\bel_t,\task),\\
  \tilde{\act}_t &= \mathrm{Shield}(\act_t,\bel_t,C_t),\\
  \ctrlcmd_t &= \controller(\tilde{\act}_t,\hat{\state}_t,C_t),\\
  y_t &= \mathrm{Evaluate}(\state_t,\act_t,\ctrlcmd_t,\task),\\
  \mathcal{D}_{k+1} &= \mathcal{D}_{k}\cup
  \mathrm{Log}(\obs_t,\bel_t,\act_t,\ctrlcmd_t,y_t).
  \end{aligned}
  \label{eq:ch21-real-vla-loop}
\end{equation}

이 loop에서 하나라도 빠지면 Real VLA claim은 약해진다. $\mathrm{Estimate}$가 없으면 perception uncertainty를 모른다. $\mathrm{Shield}$가 없으면 unsafe action을 구분하기 어렵다. $\controller$가 명시되지 않으면 action의 물리 의미가 없다. $\mathrm{Evaluate}$와 $\mathrm{Log}$가 없으면 성공과 실패가 다음 data engine으로 돌아가지 않는다.

\section{Running example revisited}

running example은 ``머그컵을 drying rack에 올리고 유리컵은 건드리지 말라''였다. 이 예제는 책 전체를 관통하는 작은 benchmark로 읽을 수 있다.

\begin{table}[h]
  \centering
  \caption{mug/rack/glass 예제를 stack 전체로 읽기}
  \begin{tabularx}{\linewidth}{p{0.22\linewidth}X X}
    \toprule
    Layer & 이 예제에서 해야 할 일 & 실패하면 보이는 현상 \\
    \midrule
    Perception & mug, rack slot, glass, free space를 찾음 & object는 찾지만 glass clearance를 모름 \\
    State/belief & pose uncertainty와 occlusion을 추정 & rack slot 위치가 불확실한데 바로 접근 \\
    Planning & pick, lift, transport, place, verify 순서를 구성 & plausible plan은 있지만 executable path가 없음 \\
    Action & delta pose, waypoint, chunk 중 interface 선택 & action quantization 또는 frame error \\
    Control & gripper, arm, contact force를 추적 & overshoot, slip, rack collision \\
    Evaluation & success, no-glass-contact, latency, recovery 기록 & 성공률은 높지만 near-miss가 숨겨짐 \\
    Safety & forbidden region과 force limit을 강제 & 유리컵을 건드리고도 success로 집계 \\
    Data engine & failure/correction/intervention을 다음 학습으로 보냄 & 같은 failure를 다음 rollout에서 반복 \\
    \bottomrule
  \end{tabularx}
  \label{tab:ch21-running-example}
\end{table}

이 예제에서 좋은 VLA는 단순히 ``mug''와 ``rack''을 이해하는 모델이 아니다. 좋은 VLA는 forbidden region, grasp approach, rack clearance, controller reference, interruption rule, safety log, evaluation metadata를 함께 남기는 system이다.

\section{Twelve questions for reading VLA papers}

VLA 논문을 읽을 때 다음 12문항을 먼저 확인한다.
\begin{enumerate}
  \item observation은 무엇이며, camera/proprioception/tactile calibration이 명시되는가?
  \item language instruction은 task goal인가, constraint인가, preference인가?
  \item state 또는 belief를 쓰는가, 아니면 raw observation만 쓰는가?
  \item action representation은 token, vector, waypoint, chunk, distribution 중 무엇인가?
  \item policy-level action이 어떤 controller command로 변환되는가?
  \item VLA inference rate와 controller rate는 각각 얼마인가?
  \item action normalization, frame convention, clipping rule이 robot-specific으로 정의되는가?
  \item data에는 success뿐 아니라 failure, correction, intervention, reset이 들어가는가?
  \item evaluation은 trial count, randomization, confidence, failure taxonomy를 제공하는가?
  \item safety는 prompt rule인가, geometric/dynamic/runtime shield인가?
  \item latency, stale action, watchdog, fallback controller가 보고되는가?
  \item deployment log가 다음 data collection 또는 adaptation으로 돌아가는가?
\end{enumerate}

이 질문 중 절반 이상이 비어 있으면, 그 논문은 VLA model claim은 할 수 있어도 Real VLA system claim은 아직 약하다.

\section{Ten propositions}

\begin{center}
  \captionof{table}{이 책의 핵심 명제 10개}
  \label{tab:ch21-ten-propositions}
  \begin{tabularx}{\linewidth}{p{8mm}X}
    \toprule
    번호 & 명제 \\
    \midrule
    1 & VLA는 더 큰 VLM이 아니라 robot execution stack에 연결되는 책임 구조이다. \\
    2 & observation, state, belief를 구분하지 않으면 perception success와 robot capability를 혼동한다. \\
    3 & $\act_t$와 $\ctrlcmd_t$는 다르며, 그 사이에는 controller contract가 필요하다. \\
    4 & contact-rich manipulation은 geometry만으로 설명되지 않는다. force, compliance, slip, recovery가 필요하다. \\
    5 & LLM/VLM plan은 executable plan이 아니다. feasibility, IK, collision, recovery가 따로 필요하다. \\
    6 & imitation learning의 covariate shift는 VLA fine-tuning에서도 핵심 failure mode이다. \\
    7 & action representation은 formatting 문제가 아니라 control boundary 문제이다. \\
    8 & robot data는 web data가 아니다. embodiment, controller, action schema, reset policy가 의미를 만든다. \\
    9 & benchmark success는 robot capability의 일부 증거일 뿐이다. metadata와 failure taxonomy가 필요하다. \\
    10 & safety는 prompt policy가 아니라 semantic, geometric, dynamic, runtime, data, human layer의 stack property이다. \\
    \bottomrule
  \end{tabularx}
\end{center}

\chapterwork
{이 장의 12 questions를 책 전체의 final paper-review protocol로 사용하라. 특히 Chapter 13, 14, 16, 19와 함께 다시 읽는 것이 좋다.}
{이 책의 10개 핵심 명제 중 실제 VLA 논문 리뷰에서 가장 자주 빠지는 명제는 무엇인가?}
{VLA 논문 하나를 골라 12 questions checklist로 short review를 작성하고, claim strength를 strong/moderate/weak으로 분류하라.}

\section{What remains}

이 장은 책의 결론이 아니라 압축이다. 지금까지 배운 것은 현재의 Real VLA system contract이다. 다음 질문은 이 contract가 앞으로 어떻게 확장되는가이다. 더 강한 world model은 어떤 state를 예측해야 하는가? Online adaptation은 어떤 release gate를 통과해야 하는가? Multi-embodiment transfer는 semantic, geometric, control transfer를 어떻게 분리해야 하는가? Humanoid VLA는 balance와 manipulation을 어떤 hierarchy로 결합해야 하는가?

따라서 다음 장의 Future Directions는 공상적 전망이 아니라, 지금까지 정리한 책임 구조 위에서 남는 연구 문제의 목록이다. 이제 우리가 배운 현재의 system contract 위에서, 앞으로 어떤 연구 문제가 남는가를 본다.
